% Generated draft: requires mathematical and source review.
\documentclass[11pt]{article}
\usepackage[margin=25mm]{geometry}
\usepackage{fontspec}
\setmainfont{Latin Modern Roman}
\usepackage{amsmath,amssymb,mathtools,unicode-math}
\setmathfont{Latin Modern Math}
\usepackage{enumitem,needspace,xurl,hyperref}
\setlength{\parindent}{0pt}
\setlength{\parskip}{5pt}
\setlength{\emergencystretch}{3em}
\newcommand{\N}{\mathbb{N}}
\newcommand{\Z}{\mathbb{Z}}
\newcommand{\Q}{\mathbb{Q}}
\newcommand{\R}{\mathbb{R}}
\newcommand{\C}{\mathbb{C}}
\newcommand{\F}{\mathbb{F}}
\newcommand{\A}{\mathbb{A}}
\newcommand{\PP}{\mathbb{P}}
\newcommand{\E}{\mathbb{E}}
\newcommand{\eps}{\varepsilon}
\newcommand{\dd}{\,\mathrm{d}}
\begin{document}
\section*{UnsolvedMath ID: 701391}
\textbf{Source problem code:} catalogue entry
\subsection{Definitions and mathematical statement}
Let \(F\) be a non-cyclic free group. Define its derived series by
\[
F^{(0)}=F,\qquad F^{(n+1)}=[F^{(n)},F^{(n)}]\quad(n\geq 0),
\]
where \([H,H]\) is the subgroup generated by the commutators of elements of \(H\). The completion intended in the source is
\[
\widehat F_{\mathrm{sol}}=\varprojlim_{n\geq 1} F/F^{(n)},
\]
with the natural quotient maps as bonding maps. These quotients are not required to be finite. For a group \(G\), its ordinary Schur multiplier is \(M(G)=H_2(G;\Z)\), computed for the underlying abstract group.

The question, under this ordinary Schur multiplier interpretation, is whether
\[
\text{for every non-cyclic free group }F,\qquad
\bigl|H_2(\widehat F_{\mathrm{sol}};\Z)\bigr|>\aleph_0.
\]
The source imposes no finite-generation or rank restriction beyond non-cyclicity. It does not explicitly discuss ordinary versus continuous homology; the displayed formulation uses ordinary group homology.

For background, define the lower central series by \(\gamma_1(F)=F\) and \(\gamma_{n+1}(F)=[\gamma_n(F),F]\). The supplied record attributes to Bousfield the assertion that the free pronilpotent completion \(\varprojlim_{n\geq 1}F/\gamma_n(F)\) of a non-cyclic free group has uncountable Schur multiplier. This is supplied background attribution, without a bibliographic citation or explicit rank qualifications for that theorem.
\subsection{Short English statement}
For every non-cyclic free group, form the inverse limit of its quotients by the derived subgroups. Must the resulting group's ordinary Schur multiplier be uncountable? The supplied record labels this question open; that status has not been independently verified.
\subsection{Sources}
\begin{itemize}
\item Local UnsolvedMath registry: permanent ID \(701391\), exact source problem code \texttt{catalogue entry}, title \emph{Kourovka Notebook Problem 17.89}. The supplied record is marked open and was updated on \texttt{2026-10-16}.
\item Supplied source attribution: \emph{The Kourovka Notebook (Unsolved Problems in Group Theory)}, original problem 17.89, 17th issue (2010); also identified with the 21st issue edition (2026), \texttt{arXiv:1401.0300v48}. Proposed by R. Mikhailov in 2010. The supplied collection description names E. I. Khukhro and V. D. Mazurov as editors.
\item The source invokes Bousfield's theorem concerning the Schur multiplier of a free pronilpotent completion, but supplies no specific bibliographic reference for that theorem. The notebook attribution and theorem attribution have not been independently checked.
\end{itemize}
\end{document}
